\scanpage{10}
$X$ surface \inserted{n’est pas une restriction essentielle ! $\dim X\geq 2$} algébrique projective non singulière conn.
sur $k$ alg. clos.

$Y_0$ section hyperplane non singulière, dans un
``pencil'' défini par $L^{n-2}\subset P^n$ coupant $X$
transversalement en $Z$

$\widetilde X$ \unsure{obtenue|obtenu} en faisant éclater $Z$

\diagram{\begin{tikzcd}[ampersand replacement=\&]
X \& \widetilde X\arrow[l,"\varphi"']\arrow[d] \& \widetilde Y_{\bar\eta}\\
\deleted{\downarrow} \& \mathbb P^1 \&
\end{tikzcd}}

\diagram{\begin{tikzcd}[ampersand replacement=\&]
\operatorname{Pic}(X)\arrow[r,hook]\arrow[dr] \& \operatorname{Pic}(\widetilde X)\arrow[d]\\
\& \operatorname{Pic}(\widetilde Y_{\overset{\deleted{-}}{\eta}})
\end{tikzcd}}

Soit $\operatorname{Pic}(X)'$ ($\operatorname{Pic}(\widetilde X)'$) le sous-groupe des
\deleted{diviseurs} classes induisant sur $\widetilde Y_{\bar\eta}$ une classe
de \uline{degré} $0$. On a donc

\diagram{\begin{tikzcd}[ampersand replacement=\&]
\operatorname{Pic}(X^{\deleted{\unsure{'|*}}})'\arrow[r,hook]\arrow[dr] \& \operatorname{Pic}(\widetilde X)'\arrow[d]\\
\& \operatorname{Pic}^0(\widetilde X_{\bar\eta})=J^0(\widetilde X_{\bar\eta})
\end{tikzcd}}
% archive sheet number: 9
